\chapter{文化知识体系构建：知识图谱的实现}\label{ux6587ux5316ux77e5ux8bc6ux4f53ux7cfbux6784ux5efaux77e5ux8bc6ux56feux8c31ux7684ux5b9eux73b0}

\section{实体与关系的定义}\label{ux5b9eux4f53ux4e0eux5173ux7cfbux7684ux5b9aux4e49}

\subsection{实体分类}\label{ux5b9eux4f53ux5206ux7c7b}

为使典籍文本中的人物、空间、技艺、文献和历史事件能够进入同一分析框架，本研究设计了一套双重分类体系。研究分类层面，实体分为六大类：非遗文化、物质遗产、传承主体、文化空间、文献记忆、历史时序，并进一步细分为 28 小类；管理口径层面，能够对应官方非遗或文物分类的实体，再补充民间文学、传统音乐、传统舞蹈、传统戏剧、曲艺、传统体育游艺与杂技、传统美术、传统技艺等官方标签。这样处理可以让典籍中的地方文化记忆既被学术分析使用，也与后续的非遗名录、文物保护单位和旅游 POI 发生连接。附录 A 给出实体分类体系简表。

\subsection{关系分类}\label{ux5173ux7cfbux5206ux7c7b}

关系分为九大语义组：

空间关联：位于、坐落于、发源于；

人物关联：出生于、任职于、交游；

传承延续：师承、传承、传授；

营建创造：始建于、创建、重建；

时序归属：形成于、兴盛于、活动于；

社群组织：属于、包含、隶属；

文献记载：记载、收录、著述;

从属分类：归入、属于、代表;

文化表征：体现、象征、承载。

关系抽取坚持开放原则，不强行把所有关系压缩为少数固定谓词，而是保留原始关系短语、关系组和证据句，以减少语义损失，并便于后续人工复核。

\section{基于大语言模型的抽取过程}\label{ux57faux4e8eux5927ux8bedux8a00ux6a21ux578bux7684ux62bdux53d6ux8fc7ux7a0b}

\subsection{模型与分块}\label{ux6a21ux578bux4e0eux5206ux5757}

实体---关系抽取使用阿里云通义千问 qwen3.5-plus 作为主要模型。由于典籍文本篇幅较长，本文将 53 份标准化语料按约 800 字一段切分，再分批送入模型处理。每一段文本进入模型后，按照预设的实体分类和关系分类输出结构化结果，内容包括实体名称、实体类别、关系短语、关系组、置信度、来源文献和原文证据句。

分块处理主要服务于两点：其一，避免长文本一次性输入造成上下文丢失；其二，让每条抽取结果保留相对明确的原文出处，便于后续校验。图 4.1 展示了从文本切块、模型抽取、质量控制到合并入库的基本流程。

\begin{figure}
\centering
\includegraphics[width=5.62083in,height=2.81042in]{ms_thesis/image3.png}
\caption{实体---关系抽取的分块与并发处理流程}
\end{figure}

\subsection{质量控制}\label{ux8d28ux91cfux63a7ux5236}

抽取过程设置了以下质量控制环节：

\begin{enumerate}
\def\labelenumi{\arabic{enumi})}
\item
  方向校验：避免主宾颠倒，如``师从 A''与``A 传于''在方向上不能混用；
\item
  禁用空泛关系：``相关''\,``有关''等无信息关系被丢弃；
\item
  双重标签：每个实体保留研究分类标签，能够对接官方口径的实体再补充官方分类标签；
\item
  证据句保留：每一条关系附带原文摘录，便于人工复核；
\item
  置信度标记：对抽取结果保留置信度信息，用于后续筛选和人工核查。
\end{enumerate}

这些控制并不能完全消除大语言模型抽取的不确定性，但能够把错误控制在可检查的范围内。对于本文而言，知识图谱不是直接替代人工判断，而是为文化记忆量化提供一套可追溯的中间层。

\subsection{合并与入库}\label{ux5408ux5e76ux4e0eux5165ux5e93}

原始抽取共得到实体记录 61,218 条、关系记录 42,024 条。由于同一人物、地点或技艺可能在不同文献中以不同写法出现，本文进行实体对齐、名称规范化、同义合并和规则筛选。合并后，知识图谱总库包含 8,048 个实体、19,382 条关系。其中，能够通过地名、建筑实体或文化锚点直接落到空间坐标的实体为 575 个。

在可视化阶段，本文按 28 个小类的语料提及频次对节点进行分级，并映射为六组高对比度颜色。节点分级和配色主要用于在全局网络中快速识别高频实体类型与关键节点，为后续的典型子图分析提供基础。

\section{图数据库的构建与可视化}\label{ux56feux6570ux636eux5e93ux7684ux6784ux5efaux4e0eux53efux89c6ux5316}

\subsection{构建过程}\label{ux6784ux5efaux8fc7ux7a0b}

图数据库构建以合并后的实体和关系总库为基础。本文将节点与关系导入 Neo4j，并完成索引建立、节点着色和层次划分。实体数量、关系数量和后续指数计算均以合并后的知识图谱总库为准；Neo4j 主要用于检查网络结构、观察局部关系和导出典型子图。

通过这一过程，典籍中的空间、人物、技艺与文献记忆被组织为一个可查询、可追溯的关系网络。它并不直接给出文旅融合结论，而是为后续``文化记忆指数''的计算提供基础材料。

\subsection{典型子图展示}\label{ux5178ux578bux5b50ux56feux5c55ux793a}

为增强图面可读性，本文将知识图谱全局网络、人物关联总图和典型人物子图分别展示。图 4.2 用于观察知识图谱整体结构和高频节点位置；图 4.3 用于说明人物谱系内部的关系组织；图 4.4 选取黄飞鸿周边人物子图作为局部网络示例，用于说明典籍记忆如何通过人物、地点、技艺与事件向外扩展。


\begin{figure}[p]
\centering
\makebox[0.96\textwidth][c]{%
\begin{minipage}[t]{0.485\textwidth}
\centering
\includegraphics[width=\linewidth]{ms_thesis/fig_4_2a_kg_global.png}
\caption{知识图谱全局网络}
\end{minipage}\hspace{0.02\textwidth}
\begin{minipage}[t]{0.485\textwidth}
\centering
\includegraphics[width=\linewidth]{ms_thesis/fig_4_2b_person_network.png}
\caption{人物关联总图}
\end{minipage}
}

\vspace{0.6em}

\centering
\includegraphics[width=0.96\textwidth]{ms_thesis/fig_4_3b_huang_feihong.png}
\caption{黄飞鸿周边人物子图}
\end{figure}

黄飞鸿子图用于解释典型文化记忆的关系来源。该子图串联洪拳、十字拳、狮舞---醒狮、广东十虎、佛山精武体育会等实体，比合并图更适合逐项识读节点标签，也更便于说明文化记忆如何从人物向地点、技艺和事件扩展。

这些子图为第 5 章的耦合分析提供了语义基础。当研究某个物质文化载体时，可以从知识图谱侧反向追溯与其关联的人物、事件、典籍与技艺实体，并将相关提及频次和关系证据转化为文化记忆指数的原始量值。知识图谱结果仍然依赖文字识别质量、模型抽取稳定性和后续规则筛选，因此本文在使用图谱指标时保留证据句和来源文献，便于人工核查。
